\documentclass[letterpaper]{article}

% AAAI-26 style. Place aaai2026.sty from the official AAAI author kit in
% this directory before compiling.
\usepackage[submission]{aaai2026}

\usepackage{times}
\usepackage{helvet}
\usepackage{courier}
\usepackage[hyphens]{url}
\usepackage{graphicx}
\urlstyle{rm}
\def\UrlFont{\rm}
\usepackage{natbib}
\usepackage{caption}
\usepackage{booktabs}
\usepackage{amsmath}
\usepackage{amssymb}
\usepackage{algorithm}
\usepackage{algorithmic}
\usepackage{newfloat}
\usepackage{listings}
\usepackage{xcolor}

\DeclareCaptionStyle{ruled}{labelfont=normalfont,labelsep=colon,strut=off}
\lstset{
  basicstyle={\footnotesize\ttfamily},
  breaklines=true,
  columns=fullflexible,
  keepspaces=true,
}

\pdfinfo{
/TemplateVersion (2026.1)
}

\setcounter{secnumdepth}{2}

\title{Consensus Is Not Completion: False Convergence in Multi-Agent Discovery Workflows}

\author{
    Anonymous Authors
}
\affiliations{
    Paper under anonymous review
}

\begin{document}

\maketitle

\begin{abstract}
Multi-agent language-model workflows often treat agreement, overlap, or a final
summary as evidence that a task is complete. This is unsafe in closed-world
discovery tasks, where the goal is to find all relevant items. We introduce
\emph{false convergence}: a workflow state in which agents report completion
and exhibit high agreement, yet the final answer omits true items. We identify
two mechanisms: \emph{aggregation loss}, where minority-discovered true items
are discarded during consensus or summarization, and \emph{common blind spots},
where agents agree because they miss the same region of the search space. On a
real Click repository audit, a standard LLM summarizer self-reports completion
in all three seeds while retaining only 68.5--74.5\% of oracle items. We argue
that completion should be modeled as coverage-risk assessment under correlated
exploration rather than vote-counting. We implement a decoupled completion
certificate with pre-audit and post-audit states: the pre-audit certificate uses
only exploration logs, while holdout gain is reserved for audit and post-audit
updates. Experiments on synthetic tasks and two real-repository audits show that
simple stopping baselines false-certify many incomplete states. The proposed
certificate avoids false certification on the current mixed-state suite, but its
safe coverage is low, exposing both the promise and the limits of the current
rule.
\end{abstract}

\section{Introduction}

Multi-agent language-model systems are increasingly used for tasks that require
search, decomposition, and synthesis. A common design pattern is simple:
multiple agents independently inspect a task, their outputs are aggregated by
majority vote or a summarizer, and the system stops when agents agree or report
high confidence. This pattern is natural. Agreement is easy to measure, and
summaries provide a clean final answer.

However, agreement is not the same as completion. In closed-world discovery
tasks, such as finding every affected file location, every policy case, or every
relevant document in a bounded corpus, the central question is not only whether
the final answer is plausible. It is whether there remains missing mass outside
the discovered set. A group of agents can converge for the wrong reason: they
may share a search strategy, omit the same boundary cases, or discard minority
evidence during aggregation.

This paper studies this failure mode. We call it \emph{false convergence}: a
workflow state in which the system appears complete under its own completion
signals but is incomplete under a closed-world oracle. The failure is subtle
because the final output can have perfect precision and high self-reported
confidence. The error is not necessarily hallucination. It is often omission.

We make three contributions.

\begin{itemize}
    \item We define false convergence in closed-world agentic discovery and
    distinguish two mechanisms: aggregation loss and common blind spots.
    \item We model completion as coverage-risk assessment rather than
    vote-counting over agent outputs.
    \item We propose a lightweight completion certificate and an
    evidence-preserving audit policy, and evaluate them on synthetic tasks and
    two real-repository audits.
\end{itemize}

Our results do not imply that union aggregation is always better than consensus.
Instead, they show that singleton evidence has two faces. Sometimes it is
missing mass; sometimes it is noise. A reliable completion controller should
model this distinction rather than silently deleting minority evidence or
blindly accepting all unique items.

\section{Related Work}

\subsection{Multi-Agent Consensus, Aggregation, and Calibration}

Self-consistency selects a consistent answer from multiple samples
\citep{wang2022selfconsistency}; debate, collaboration, and role-specialized
agent frameworks improve reasoning or task performance
\citep{du2023debate,chan2023chateval,wu2023autogen,chen2023agentverse,
hong2023metagpt,li2023camel,zhang2024proagent}. Recent work also studies
majority-failure auditing, selection bottlenecks, trace-level synthesis, and
calibration under communication-induced correlation
\citep{yang2026agentauditor,maryanskyy2026agentsdisagree,
fadnavis2026beyondconsensus,huang2026cagecal}. Tool-using and
planning agents extend LLMs through environment interaction
\citep{yao2023react,schick2023toolformer,qin2023toolllm,yao2023tot,
wang2023voyager}. Benchmarks such as
AgentBench, WebArena, GAIA, and SWE-bench stress interactive or software tasks
\citep{liu2023agentbench,zhou2023webarena,mialon2023gaia,jimenez2023swebench}.
Our focus is different: aggregation is judged by closed-world recall and
calibrated stopping risk, not by the quality of one selected answer.

\subsection{Wide Research and Completeness Evaluation}

Self-Refine, Reflexion, CRITIC, Chain-of-Verification, and related factuality
checks repair or verify outputs
\citep{madaan2023selfrefine,shinn2023reflexion,gou2023critic,dhuliawala2023cov,
lin2021truthfulqa,manakul2023selfcheckgpt,li2023halueval}. Retrieval-augmented
generation and dense retrieval reduce hallucination but still face noise,
negative rejection, integration, and long-context failures
\citep{lewis2020rag,guu2020realm,karpukhin2020dpr,nakano2021webgpt,
chen2024rgb,li2024unigen,liu2023lost,asai2023selfrag,
gao2023ragsurvey}. Closed-world discovery is also related to high-recall
retrieval, fact verification, and technology-assisted review, where stopping is
a coverage-estimation problem
\citep{thorne2018fever,yang2018hotpotqa,grossman2015trec,
cormack2021tarstopping}. Our certificate adopts this stopping perspective for
correlated LLM-agent exploration.
Recent broad or deep research benchmarks explicitly evaluate exhaustive
information seeking and stopping uncertainty
\citep{wong2025widesearch,gupta2026deepsearchqa,rafiee2026trqa,
huang2026wideseek,kim2026seekergym}.

\subsection{Coverage Estimation and Stopping}

Our certificate uses observable missing-mass signals rather than oracle labels at
decision time. Singleton counts and unseen-mass estimates are conceptually
related to Good-Turing, Chao-style, capture-recapture, and ecological diversity
estimators \citep{good1953population,chao1984nonparametric,
mcAllester2000goodturing,chao2014rarefaction,bron2024chao}. Recent LLM and
multi-agent uncertainty methods
estimate answer risk from probabilities, samples, or semantic variation
\citep{kuhn2023semanticentropy,kadavath2022language,chen2026matu,
cole2023selectively}. We study a complementary question: whether a multi-agent
workflow has covered a bounded item set well enough to stop.

\section{Problem Formulation}

\subsection{Closed-World Discovery Tasks}

We consider tasks with a bounded source collection $D$ and a finite oracle set
$Y^\star$ of target items. A run returns a set $\hat{Y}$. We evaluate:

\begin{align}
\text{Recall}(\hat{Y}) &= \frac{|\hat{Y} \cap Y^\star|}{|Y^\star|}, \\
\text{Precision}(\hat{Y}) &= \frac{|\hat{Y} \cap Y^\star|}{|\hat{Y}|}.
\end{align}

Examples include code dependency audits, policy-case discovery, high-recall
document review, and bounded repository analysis. Unlike open-ended question
answering, a plausible final answer is insufficient: the system must decide
whether it has found the complete set.

\subsection{Multi-Agent Completion Signals}

A typical workflow runs $k$ agents $A_1, \ldots, A_k$, each producing an item set
$Y_i$, a completion flag $c_i$, and a confidence score $q_i$. Aggregation then
produces a final set $Y_{\text{agg}}$. Common aggregation rules include:

\begin{itemize}
    \item \textbf{Majority consensus}: keep items reported by at least two
    agents.
    \item \textbf{Standard summarization}: ask a summarizer to produce a clean
    final answer from the agent reports.
    \item \textbf{Raw union}: keep every unique item reported by any agent.
    \item \textbf{Union-preserving summarization}: instruct a summarizer to
    preserve all unique reported items unless malformed or contradicted.
\end{itemize}

Systems often treat high pairwise overlap, high mean confidence, and affirmative
self-reported completion as stopping evidence. We argue that these are not
completion certificates.

\subsection{False Convergence}

Let $S$ be a multi-agent workflow state with $k$ agents, an aggregation rule
$g$, and a target recall threshold $\theta$. We say $S$ exhibits false
convergence under $g$ when:

\begin{enumerate}
    \item agents self-report completion with high confidence;
    \item agent outputs exhibit high overlap;
    \item the aggregated answer $g(Y_1,\ldots,Y_k)$ has recall below $\theta$.
\end{enumerate}

Our preregistered scoring uses $\theta=0.95$, confidence threshold $0.80$, and
overlap threshold $\tau=0.70$. Holdout recovery is a diagnostic and audit
mechanism, not part of the definition. We also report near-positive and
mechanism-positive cases when the missing items are few but diagnostically
important.

\subsection{Aggregation Loss}

Aggregation loss occurs when a true item appears in at least one agent report
but is removed by consensus-style aggregation or standard summarization. This is
not a search failure: the evidence exists in the reports. The failure happens at
the final synthesis step.

Formally, an oracle item $y \in Y^\star$ is an aggregation-loss item if:
\begin{equation}
y \in \bigcup_i Y_i \quad \text{and} \quad y \notin Y_{\text{agg}}.
\end{equation}
When $y$ appears in only one agent report, it is a singleton true positive.

\subsection{Common Blind Spots}

Common blind spots occur when multiple agents agree because they share a search
failure. In this case, the missing item is absent even from the raw union:
\begin{equation}
y \in Y^\star \quad \text{and} \quad y \notin \bigcup_i Y_i.
\end{equation}
Raw union and union-preserving summarization cannot recover such items because
there is no reported evidence to preserve. The mitigation must change the search
or auditing process, for example through boundary-focused holdout review.

\subsection{Singletons as Missing Mass or Noise}

Singletons are ambiguous. A singleton may be a hard-to-find true positive, or it
may be a false positive. Consensus implicitly treats singletons as noise. Raw
union implicitly treats them as valid evidence. Both assumptions are unsafe. The
right action is often to preserve singleton evidence for audit.

\section{Coverage-Risk Completion Certificate}

We propose a lightweight completion certificate for deciding whether a
multi-agent discovery workflow can safely stop. The certificate is not an
oracle: it does not use ground truth and does not decide final item labels.
Instead, it maps observable workflow signals to a stopping judgment.

\begin{figure*}[t]
\centering
\setlength{\unitlength}{0.7pt}
\begin{picture}(720,125)
\put(5,78){\framebox(105,34){Agent reports}}
\put(135,96){\vector(1,0){60}}
\put(205,78){\framebox(125,34){Aggregation loss}}
\put(205,38){\framebox(125,34){Common blind spot}}
\put(350,96){\vector(1,0){55}}
\put(350,56){\vector(1,0){55}}
\put(415,78){\framebox(105,34){Evidence ledger}}
\put(535,96){\vector(1,0){45}}
\put(590,78){\framebox(120,34){Pre-audit cert.}}
\put(590,38){\framebox(120,34){Audit policy}}
\put(650,78){\vector(0,-1){20}}
\put(650,38){\vector(0,-1){20}}
\put(550,5){\framebox(160,28){Post-audit certificate}}
\put(62,63){\makebox(0,0){minority true item dropped}}
\put(62,23){\makebox(0,0){shared missing region}}
\end{picture}
\caption{False convergence mechanisms and the decoupled completion workflow.
Holdout gain is not a pre-audit signal; it updates the ledger only after audit.}
\label{fig:workflow}
\end{figure*}

\subsection{Observable Signals}

The pre-audit certificate takes as input only exploration-stage reports
$Y_1,\ldots,Y_k$, confidence values, and source metadata. It does not use oracle
labels or holdout gain. We compute item counts $n(y)$ and frequency counts
$f_r=|\{y:n(y)=r\}|$. The Chao-style unseen estimate is
\begin{equation}
\hat{U}_{\text{Chao}} =
\begin{cases}
f_1^2/(2f_2), & f_2>0,\\
f_1(f_1-1)/2, & f_2=0,
\end{cases}
\end{equation}
with missing ratio
$m_{\text{Chao}}=\hat{U}_{\text{Chao}}/(|\cup_iY_i|+\hat{U}_{\text{Chao}})$.
We also compute Good--Turing missing mass $m_{\text{GT}}=f_1/\sum_y n(y)$,
mean output Jaccard, mean source overlap, source coverage, and an effective
exploration size $k_{\text{eff}}=k/(1+(k-1)\rho)$, where $\rho$ is mean output
Jaccard.

\subsection{Certificate Rule}

The implemented rule computes a continuous risk score
\begin{align}
R = &\ r_{\text{conf}} + .25\min(1,s/0.20)
      + .22\min(1,m_{\text{Chao}}/0.10) \nonumber\\
    &+ .20\min(1,m_{\text{corr}}/0.15)
      + r_{\text{src}} + r_{\text{boundary}} + r_{\text{audit}},
\end{align}
clipped to $[0,1]$, where $s=f_1/|\cup_iY_i|$ and
$m_{\text{corr}}=\min(1,m_{\text{Chao}}\cdot k/k_{\text{eff}})$. The confidence
penalty $r_{\text{conf}}$ is 0.12 when mean confidence is below 0.80. The source
penalty is proportional to missing source coverage below 0.30. Boundary risk
adds 0.12 when mean overlap is at least 0.95 on a boundary-sensitive task.
Post-audit unresolved singleton evidence adds 0.12.

The label is \textsc{safe-to-stop} only when confidence is high,
$m_{\text{corr}}\le0.05$, singleton evidence is absent or resolved by audit, and
$R<0.25$. It is \textsc{unsafe-to-stop} when $R\ge0.60$ or
$m_{\text{corr}}\ge0.20$; otherwise it returns \textsc{requires-audit}. These
thresholds are fixed in code and written to the result log.

\subsection{Certificate Outputs}

The certificate returns one of three labels:

\begin{itemize}
    \item \textsc{safe-to-stop}: observable risk is low enough that the workflow
    may publish the current conservative answer.
    \item \textsc{unsafe-to-stop}: observable risk suggests likely missing
    mass; the workflow should continue searching or run a stronger audit.
    \item \textsc{requires-audit}: the current answer may be complete, but
    singleton evidence, noisy union items, or weak independence prevents clean
    certification.
\end{itemize}

In the current experiments, the certificate is deliberately conservative. It is
designed to catch high-risk states rather than to maximize automatic stopping.
We therefore report false-certification rate and safe coverage separately.

\subsection{Evidence-Preserving Audit Policy}

When the certificate returns \textsc{unsafe-to-stop} or
\textsc{requires-audit}, the workflow needs an audit policy. We use an
evidence-preserving policy that separates finalization from evidence retention:
consensus items form a conservative answer, singleton evidence is retained for
review, and high-agreement boundary-risk states trigger holdout or source-focused
audit. This policy is the audit response to a risky certificate state, not the
certificate itself.

\begin{algorithm}[H]
\caption{Evidence-Preserving Audit Policy}
\begin{algorithmic}[1]
\STATE Run $k$ independent discovery agents, producing $Y_1,\ldots,Y_k$.
\STATE Compute item counts $n(y)=|\{i:y\in Y_i\}|$.
\STATE Set conservative final set $F=\{y:n(y)\ge2\}$.
\STATE Set audit queue $Q=\{y:n(y)=1\}$.
\STATE Compute pre-audit confidence, overlap, source coverage, and unseen-mass signals.
\STATE Apply the pre-audit completion certificate.
\IF{certificate is \textsc{safe-to-stop}}
    \STATE Return $F$ with certificate report.
\ENDIF
\IF{$Q$ is non-empty}
    \STATE Trigger singleton audit.
\ENDIF
\IF{confidence and overlap are high and the task has boundary-resolution risk}
    \STATE Trigger boundary-focused holdout audit.
\ENDIF
\STATE Add audit-verified singletons to $F$.
\STATE Add holdout-discovered boundary items to $F$.
\STATE Recompute post-audit certificate using the updated ledger.
\STATE Return $F$ plus pre/post certificate reports.
\end{algorithmic}
\end{algorithm}

The audit policy handles the two failure mechanisms differently. For aggregation
loss, it preserves and audits singleton or minority evidence instead of deleting
it during summarization. For common blind spots, where no current agent reported
the missing item, it triggers boundary-focused holdout or source-focused review.
In our experiments, singletons are not automatically accepted; they are retained
and audited. High agreement is not automatically trusted; it can trigger a
common-blind-spot check.

\section{Experimental Setup}

\subsection{Research Questions}

We ask:

\begin{itemize}
    \item \textbf{RQ1}: Can consensus-style aggregation produce false completion
    signals in closed-world discovery tasks?
    \item \textbf{RQ2}: Are failures caused by aggregation loss, common blind
    spots, or both?
    \item \textbf{RQ3}: Does raw union solve the problem, or does it introduce a
    precision cost?
    \item \textbf{RQ4}: Can a coverage-risk certificate identify unsafe
    stopping states, and when can an audit policy recover missing items?
\end{itemize}

\subsection{Tasks}

We use four bounded discovery task families: two synthetic families designed to
isolate mechanisms, and two real-repository audits used to test whether the
failure appears beyond constructed examples.

\textbf{T1: Code dependency audit.} Agents inspect a synthetic AcmePay migration
repository and must find every code or configuration location that resolves to
AcmePay v1 through direct calls, routing configuration, registry indirection,
fallback queues, or replay scripts. The hard oracle has 35 items.

\textbf{T2: Policy document search.} Agents inspect naturalistic incident notes,
changelogs, appendices, and policy memos. They must find all cases whose state,
flow, and service-resolution chain imply AcmePay v1 handling. The oracle has 30
items. Boundary cases require following service catalog and adapter registry
resolution rather than relying on surface service names.

\textbf{T3: Real-repository Click deprecation audit.} Agents inspect a fixed
snapshot of the public Click repository and must find line-level locations that
belong to a deprecated API surface audit. The oracle has 149 items spanning
implementation, warning emission, help-message formatting, parameter behavior,
tests, documentation, and bounded changelog entries. Unlike T1/T2, T3 uses a
real repository and is used as a stronger validation of aggregation behavior
rather than as a strict false-convergence construction.

\textbf{T4: Real-repository Requests TLS audit.} We additionally instantiate a
second real-repository task family from a fixed snapshot of the public Requests
repository. The task asks agents to find line-level items related to TLS
certificate verification, CA bundle resolution, client certificates, SSL error
handling, tests, test infrastructure, and user documentation. The oracle has
304 items. T4 is used to test whether the real-repository result depends on
Click or deprecation-surface search.

\subsection{Runs and Aggregation Policies}

Each seed contains three G3 agents. Where available, we also run an independent
G6 holdout scout. To test whether the certificate simply always refuses to stop,
we construct mixed evaluation states from saved logs: each single agent, each
two-agent union, the three-agent pre-audit union, and the post-holdout state
when available. This yields complete, near-complete, and incomplete states
without generating new oracle-dependent agent outputs. We compare:

\begin{itemize}
    \item majority consensus;
    \item blind standard summarization;
    \item raw union;
    \item blind union-preserving summarization;
    \item holdout scout;
    \item self-reported completion, confidence-only, overlap-only, no-new-item,
    Chao-only, always-unsafe, always-holdout, and the proposed certificate.
\end{itemize}

The standard and union-preserving summarizers are blind to oracle and scoring
files. The pre-audit certificate uses only exploration outputs and metadata;
holdout gain is used only after audit. Oracle is used only for offline metrics.

\section{Results}

\subsection{Mechanism Evidence}

The six mechanism slices show the two failure modes. T1 seed01 and
T2 seed01 show aggregation loss: true items were present in minority reports but
dropped by consensus and standard summaries. T2 seed02 and seed03 show common
blind spots: all three G3 agents agreed on the same incomplete 28-item set, and
raw union could not recover the missing cases. T1 seed02 and seed03 are
precision-cost controls: consensus is complete, while raw union preserves false
positive singletons.

\subsection{Aggregation Policy Comparison}

Table~\ref{tab:methods} shows the method-level comparison. Consensus and
standard summarization can produce high-precision false stops. Raw union and
union-preserving summarization recover aggregation-loss cases, but incur
precision cost in T1 seed02/03 and cannot recover T2 common blind spots. The
evidence-preserving audit policy recovers missing items when holdout evidence is
available and marks singleton-noise controls as requiring audit rather than
silently accepting false positives.

\begin{table*}[t]
\centering
\small
\begin{tabular}{llrrrrl}
\toprule
Case & Method & Found & TP & FP & Recall & Status \\
\midrule
T1 seed01 & consensus / standard & 22 & 22 & 0 & 0.629 & false stop \\
T1 seed01 & raw union / union-preserving & 35 & 35 & 0 & 1.000 & complete \\
T1 seed01 & audit policy & 35 & 35 & 0 & 1.000 & verified \\
\midrule
T1 seed02 & consensus / standard & 35 & 35 & 0 & 1.000 & complete \\
T1 seed02 & raw union / union-preserving & 39 & 35 & 4 & 1.000 & precision cost \\
T1 seed02 & audit policy & 35 & 35 & 0 & 1.000 & requires audit \\
\midrule
T1 seed03 & consensus / standard & 35 & 35 & 0 & 1.000 & complete \\
T1 seed03 & raw union / union-preserving & 40 & 35 & 5 & 1.000 & precision cost \\
T1 seed03 & audit policy & 35 & 35 & 0 & 1.000 & requires audit \\
\midrule
T2 seed01 & consensus / standard & 28 & 28 & 0 & 0.933 & false stop \\
T2 seed01 & raw union / union-preserving & 30 & 30 & 0 & 1.000 & complete \\
T2 seed01 & audit policy & 30 & 30 & 0 & 1.000 & verified \\
\midrule
T2 seed02 & consensus / raw union & 28 & 28 & 0 & 0.933 & common blind spot \\
T2 seed02 & holdout / audit policy & 30 & 30 & 0 & 1.000 & verified \\
\midrule
T2 seed03 & consensus / raw union & 28 & 28 & 0 & 0.933 & common blind spot \\
T2 seed03 & holdout / audit policy & 30 & 30 & 0 & 1.000 & verified \\
\bottomrule
\end{tabular}
\caption{Method comparison on current case slices. Grouped rows combine methods
with identical scores in the current experiments.}
\label{tab:methods}
\end{table*}

\subsection{Real-Repository Summarizer False Stops}

T3 evaluates whether the aggregation-stage failure appears in a real repository
task rather than only in constructed mechanism slices. We give the same G3
aggregation packet to two blind LLM summarizers. The \emph{standard} summarizer
is instructed to produce a concise, high-precision final answer from the three
agent reports. The \emph{union-preserving} summarizer is instructed to preserve
all unique reported items unless malformed or contradicted.

Table~\ref{tab:t3-summarizers} shows a stable pattern across three blind seeds.
The standard summarizer reports completion with confidence between 0.93 and
0.96, but its recall is only 0.685--0.745. Its precision is nearly perfect, so
the failure is not hallucination; it is omission. The union-preserving
summarizer removes the false stop in all three seeds, reaching 0.993--1.000
recall, but precision falls to 0.639--0.744. This supports two claims. First,
much of the missing evidence was already present in the agent reports and was
discarded during aggregation. Second, naive union preservation is diagnostic but
not a complete solution because it preserves noise along with missing true
items.

\begin{table*}[t]
\centering
\small
\begin{tabular}{llrrrrrrr}
\toprule
Seed & Policy & Conf. & Found & TP & FP & Recall & Precision & False Stop \\
\midrule
01 & standard summarizer & 0.930 & 105 & 104 & 1 & 0.698 & 0.990 & yes \\
01 & union-preserving summarizer & 0.910 & 199 & 148 & 51 & 0.993 & 0.744 & no \\
02 & standard summarizer & 0.960 & 111 & 111 & 0 & 0.745 & 1.000 & yes \\
02 & union-preserving summarizer & 0.950 & 233 & 149 & 84 & 1.000 & 0.639 & no \\
03 & standard summarizer & 0.930 & 102 & 102 & 0 & 0.685 & 1.000 & yes \\
03 & union-preserving summarizer & 0.840 & 232 & 149 & 83 & 1.000 & 0.642 & no \\
\midrule
Mean & standard summarizer & -- & -- & -- & -- & 0.709 & 0.997 & 3/3 \\
Mean & union-preserving summarizer & -- & -- & -- & -- & 0.998 & 0.675 & 0/3 \\
\bottomrule
\end{tabular}
\caption{T3 real-repository summarizer comparison. The oracle contains 149
line-level Click deprecation-audit items. Standard summarization is
high-precision but repeatedly false-stops; union-preserving summarization
recovers recall but incurs a substantial precision cost.}
\label{tab:t3-summarizers}
\end{table*}

\subsection{Second Real-Repository Task Family}

T4 evaluates a different real repository and audit surface: TLS and certificate
handling in Requests. Unlike T3, this task is larger and noisier. Table
~\ref{tab:t4-realrepo} shows that the failure shifts from mostly aggregation
loss to mixed search-coverage and aggregation risk. Majority consensus is
incomplete in all three seeds, with recall between 0.678 and 0.701. Raw union
improves recall to 0.829--0.859, showing that some omitted evidence is present
in minority reports, but raw union remains below the 0.95 completion threshold
in every seed. Thus T4 contains candidate-pool missing mass that union
aggregation cannot recover.

The standard summarizer also false-stops in all three seeds, with recall
0.368--0.648 despite self-reported completion. The completion certificate marks
all three seeds \textsc{unsafe-to-stop}. This supports the broader claim that a
completion controller must estimate coverage risk, not merely summarize or vote
over the current reports. T4 should not be read as a direct copy of the T3
mechanism; instead, it is a harder boundary case showing that real-repository
discovery can combine aggregation loss, noisy singleton evidence, and
insufficient search coverage.

\begin{table*}[t]
\centering
\small
\begin{tabular}{lrrrrrr}
\toprule
Seed & Cons. Rec. & Std. Rec. & Union Rec. & Union Prec. & Holdout Rec. & Cert. \\
\midrule
01 & 0.701 & 0.368 & 0.859 & 0.661 & 0.750 & unsafe \\
02 & 0.678 & 0.622 & 0.845 & 0.668 & 0.688 & unsafe \\
03 & 0.678 & 0.648 & 0.829 & 0.676 & 0.763 & unsafe \\
\midrule
Mean & 0.685 & 0.546 & 0.844 & 0.668 & 0.734 & 3/3 unsafe \\
\bottomrule
\end{tabular}
\caption{T4 Requests TLS/certificate audit on a second real repository. Raw
union improves over consensus but remains incomplete, indicating missing mass
outside the G3 candidate pool. Standard summarization false-stops in all three
seeds.}
\label{tab:t4-realrepo}
\end{table*}

\subsection{Decoupled Certificate and Stopping Baselines}

We then evaluate the decoupled v1 certificate on 96 saved states constructed
from existing runs: 86 pre-audit states and 10 post-audit states. Oracle labels
are used only for evaluation; a state is safe when recall is at least 0.95. The
suite contains 44 safe and 52 unsafe states. Table~\ref{tab:v1-certificate}
shows that simple stopping signals false-certify many incomplete states:
self-reported completion and confidence-only stopping each false-certify 35
states, Chao-only false-certifies 13, and overlap-only false-certifies 10. The
proposed certificate certifies one post-audit state and makes no false
certifications, but its safe coverage is only 0.023. Thus the current rule is a
conservative risk detector, not a solved stopping policy. Its continuous risk
score has AUROC 0.515 and AUPRC 0.589 for unsafe-state detection, indicating
that stronger correlation and coverage features are still needed.

\begin{table}[t]
\centering
\small
\begin{tabular}{lrrrr}
\toprule
Rule & Stop & False & FCR & Safe Cov. \\
\midrule
Self-report & 70 & 35 & .500 & .795 \\
Confidence & 70 & 35 & .500 & .795 \\
Overlap & 15 & 10 & .667 & .114 \\
No-new-item & 22 & 10 & .455 & .273 \\
Chao-only & 26 & 13 & .500 & .295 \\
Raw union & 86 & 49 & .570 & .841 \\
Always unsafe & 0 & 0 & .000 & .000 \\
Always holdout & 10 & 3 & .300 & .159 \\
Ours v1 & 1 & 0 & .000 & .023 \\
\bottomrule
\end{tabular}
\caption{Stopping evaluation on 96 mixed saved states. FCR is false
certification rate among stopped/certified states. The proposed certificate is
safe but too conservative in this version.}
\label{tab:v1-certificate}
\end{table}

\subsection{Completion Certificate on T3}

The real-repository results also reveal a limitation of audit-by-holdout alone.
In T3 seed01, the evidence-preserving audit policy reaches 0.993 recall and 0.955
precision by auditing singleton evidence. In seed02 and seed03, however, the
available holdout scout is itself incomplete, so the audit policy recovers only to
0.805 and 0.832 recall. The completion certificate catches this boundary: for
all three T3 seeds it returns \textsc{unsafe-to-stop}, triggered by low mean
confidence, singleton missing-mass signals, and Chao unseen-mass estimates.
Thus the method contribution should be read in two layers: evidence preservation
can recover aggregation-stage omissions when audit evidence is strong, while the
certificate prevents high-risk states from being falsely certified as complete.

\subsection{Source-Aware Audit Prototype}

The T3 certificate tells the workflow not to stop, but a practical system also
needs an audit policy that converts retained evidence into recovery. We evaluate
an offline source-aware audit prototype over G3/holdout candidates, using
deterministic task-policy predicates over source lines. This is an audit-policy
upper bound, not a blind LLM result. On T3, the raw candidate pools average 0.998
recall but only 0.660 precision. Source-aware filtering keeps the same average
recall and raises precision to 1.000; a bounded source sweep reaches 1.000 recall
and 1.000 precision. On T4, candidate filtering reaches 0.849 recall and 1.000
precision, showing that precise audit is possible but search coverage remains
insufficient.

\subsection{Audit Policy Behavior}

The audit policy exposes rather than hides uncertainty. In T1 seed01, it audits and
verifies 13 singleton true positives. In T2 seed01, it verifies 2 singleton true
positives. In T2 seed02/03, there are no singletons, but high agreement and high
confidence trigger boundary-focused holdout, recovering 2 new items in each
seed. In T1 seed02/03, the audit policy leaves 4 and 5 singleton items in the audit
queue and does not mark completion, avoiding raw union's false positives.

\subsection{Ablation and Cost Proxy}

We additionally run audit-policy ablations. Removing singleton audit leaves the
aggregation-loss singletons in T1 seed01 and T2 seed01 unrecovered. Removing the
common-blind-spot trigger leaves T2 seed02/03 at 28/30 recall despite perfect G3
agreement. Removing holdout converts hidden risk into a \textsc{requires-audit}
state but cannot recover missing items. Thus the two triggers correspond to
distinct failure mechanisms, and holdout is the mechanism that turns retained
risk into verified recovery.

Current logs do not contain reliable token or wall-clock accounting, so we
report a proxy cost: audit queue size plus one unit per triggered holdout. Under
this proxy, T1 seed01 uses 14 audit actions to recover 13 true positives, T2
seed01 uses 3 actions to recover 2 true positives, and T2 seed02/03 use one
holdout unit to recover 2 true positives each. In T1 seed02/03, the audit policy
avoids the 4 and 5 false positives introduced by raw union, but correctly
refuses to certify completion without further audit.

\section{Discussion and Limitations}

\subsection{Why Consensus Fails}

Consensus can fail for two opposite reasons. It can be too conservative when a
minority report contains real evidence. It can also be overconfident when agents
are correlated and share a blind spot. These are not contradictions. They show
that agreement is evidence of shared sampling behavior, not necessarily evidence
of low missing mass.

\subsection{Why Raw Union Is Insufficient}

Raw union is a useful diagnostic: if raw union recovers missing items, then the
failure occurred after discovery, at aggregation. But raw union is not a general
solution. It can preserve singleton false positives, and it cannot recover items
that no agent reported. A production workflow needs an audit policy, not merely
a larger set operation.

\subsection{Completion as Coverage-Risk Assessment}

The core lesson is that completion should be represented as a risk assessment.
A final answer should be accompanied by signals such as singleton count,
singleton audit status, pairwise overlap, source coverage, and known
boundary-risk triggers. Holdout gain is a post-audit signal, not a pre-audit
certificate input. This reframes agent completion from a consensus problem to an
uncertainty-estimation problem under correlated exploration.

\subsection{Limitations and Next Experiments}

The evidence supports mechanism-level claims, real-repository false stops, and a
proof-of-concept coverage-risk certificate, not a general proof of safe
completion. Broader claims require more repositories, task families, seeds,
agent counts, prompt-diverse and model-heterogeneous runs, blind
source-partitioned audits, systematic token and wall-clock accounting, stronger
stopping baselines, and correlation metrics beyond output Jaccard.
We prepared a SeekerGym adapter scaffold but did not run the benchmark, so no
public-benchmark results are claimed here.

\section{Conclusion}

Consensus is not completion. In closed-world discovery tasks, multi-agent
agreement and high confidence can hide missing true items. We identified two
mechanisms: aggregation loss, where true singleton evidence is discarded, and
common blind spots, where agents agree because they miss the same boundary
cases. We showed that raw union diagnoses some aggregation failures but can
preserve false positives and cannot recover unreported items. A coverage-risk
certificate separates the stopping decision from the audit policy: it can mark
high-risk states as unsafe or requiring audit, while evidence-preserving audit
retains minority evidence and triggers additional review. These results suggest
that reliable agentic workflows need completion certificates grounded in
coverage risk, not only consensus.

\bibliography{references}
\end{document}
